% !TeX root = GeoDynamics.tex
\section{Multi-Particle Contacts}

Our model states that the dynamics of a particle depends only on the area of overlap into another particle. 


Momentum is summed The momentum in collision must equal the amount total momentum per area in rebound. But in multi-contact, a source particle cannot independently apply the full area rule to every contact as if each contact owns the whole source response capacity. That double-counts available momentum. For a particle with multiple active contacts, the source has one current momentum state and multiple overlap areas. The contacts must share the source response by weight.  This is the fundamental source of conservation of the whole system. Therefore, schedule wise, \textbf{source owned contact resolution must work}. In multiple particle contacts the weighing of area, changes the momentum per area.

In the geometric compression model, the momentum at a given frame is computed from the current overlap area, not by integrating every prior step.So it is instantaious. In the geometric compression model, the momentum at a given frame is computed from the current overlap area, not by integrating every prior step.
 

But, 

Currently multi-particle contacts averages the velocity of each particle pair.    
This is not a geometric solution but a velocity prediction. 

For a contact, the change in normal velocity is proportional to the penetration area. 
The momentum response associated with that contact is proportional to mass times that velocity change.
The relative contribution of each contact is proportional to its overlap area divided by the total active overlap area for the source particle.

The total overlap area is the sum of all target particle
overlap areas. From this the total momentum imposed on the target particle can be found.
Taking the ratio of total overlap area to the individual overlap areas gives the fraction
each target particle contributes to the collision.

From section \ref{Fundamental}, we have three relationships:
\begin{enumerate}
\item Velocity/area relation
\item Velocity change proportional to penetration area
\item Force/penetration law
\end{enumerate}

\subsection{Velocity}

The clearest rule for Mutli-Particle Contact is from the second fundamental relation:
\begin{equation}
\begin{aligned}
(v_{in} - v_p) = P A_P
\end{aligned}
\end{equation}

so that :
\begin{equation}
\begin{aligned}
\Delta v = P A_P
\end{aligned}
\end{equation}

That gives the clean area relationship:

\textbf{Velocity change is proportional to penetration/overlap area}

\subsection{Momentum}

For momentum:

\begin{equation}
\begin{aligned}
\Delta p = m \Delta v
\end{aligned}
\end{equation}

so that:


\begin{equation}
\begin{aligned}
\Delta p = m P A_P
\end{aligned}
\end{equation}

So the momentum relationship states:

\textbf{Contact momentum contribution is proportional to overlap area}

For several contacts on the same source particle:

\begin{equation}
\begin{aligned}
A_T = \sum_k A_k
\end{aligned}
\end{equation}

Where each contact’s geometric weight is:
\begin{equation}
\begin{aligned}
w_k = \frac{A_k}{A_T}
\end{aligned}
\end{equation}

Which states that:

\textbf{Each target contact contributes a fraction of the source particle's total contact response according to its overlap-area weight.}

From the local relation:
\begin{equation}
\begin{aligned}
\label{none}
\Delta v = P A_P
\end{aligned}
\end{equation}

and since:
\begin{equation}
\begin{aligned}
\label{none}
\Delta p = m \Delta v
\end{aligned}
\end{equation}

we get:


\begin{equation}
\begin{aligned}
\label{none}
\Delta p = m P A_P
\end{aligned}
\end{equation}

So the simplest area-to-momentum relationship is:

\begin{equation}
\begin{aligned}
\label{none}
\Delta p \propto m A_P
\end{aligned}
\end{equation}

But for contact dynamics, the better mass term is probably not just source mass. For a pair contact, the normal exchange depends on reduced mass:

\begin{equation}
\begin{aligned}
\label{none}
\mu_{ij} = \frac{m_i m_j}{m_i + m_j}
\end{aligned}
\end{equation}

So for a contact $k$ between source particle $i$ and target particle $j$:

\begin{equation}
\begin{aligned}
\label{none}
\Delta p_k = C \mu_{ij} A_k
\end{aligned}
\end{equation}

where:

C = contact stiffness / momentum-per-area proportionality
$A_k$ = overlap area for that contact
Then for all contacts owned by source particle $i$:

\begin{equation}
\begin{aligned}
\label{none}
\Delta p_{total,i}= C \sum_k \mu_{ik} A_{ik}
\end{aligned}
\end{equation}

That gives a geometric total contact momentum scale.

Then each contact gets:

\begin{equation}
\begin{aligned}
\label{none}
w_k =
\frac{\mu_{ik} A_{ik}}
{\sum_j \mu_{ij} A_{ij}}
\end{aligned}
\end{equation}

and:

\begin{equation}
\begin{aligned}
\label{none}
\Delta p_k = w_k \Delta p_{total,i}
\end{aligned}
\end{equation}

But because $w_k$ is built from the same quantity, this simplifies to:


\begin{equation}
\begin{aligned}
\label{none}
\Delta p_k = C \mu_{ik} A_{ik}
\end{aligned}
\end{equation}

The total is just the sum:

\begin{equation}
\begin{aligned}
\label{none}
\Delta p_{total,i}=\sum_k \Delta p_k
\end{aligned}
\end{equation}

So the real decision is the coefficient $C$.

In the old code, the closest thing is:
\begin{verbatim}
overlap_contact_momentum =
    momentum_per_area * overlap_area * inverse_square_weight(center_distance)
\end{verbatim}
So the current practical relationship is:

\begin{equation}
\begin{aligned}
\label{none}
\Delta p_k=P A_k W(d_k)
\end{aligned}
\end{equation}

where:

\begin{equation}
\begin{aligned}
\label{none}
W(d_k) = \frac{1}{\max(d_k^2, \epsilon^2)}
\end{aligned}
\end{equation}

But if we want to move more geometric and less heuristic, the cleaner relationship is:

\begin{equation}
\begin{aligned}
\label{none}
\Delta p_k = P A_k
\end{aligned}
\end{equation}

or with reduced mass:

\begin{equation}
\begin{aligned}
\label{none}
\Delta p_k = P \mu_{ik} A_k
\end{aligned}
\end{equation}

For equal masses, reduced mass is constant, so the weighting reduces to area weighting.

The important correction for multi-particle contact is:

text


Do not average pair velocity predictions.
Compute a contact momentum vector from area for each active contact,
sum those contact momentum vectors,
then apply the resulting source momentum change once.
Vector form:


\begin{equation}
\begin{aligned}
\label{none}
\Delta \mathbf{p}_{i}=-\sum_k\left(P \mu_{ik} A_{ik}\right)\mathbf{n}_{ik}
\end{aligned}
\end{equation}

where $\mathbf{n}_{ik}$ points from source to target.

Then:

\begin{equation}
\begin{aligned}
\label{none}
\Delta \mathbf{v}_{i}=\frac{\Delta \mathbf{p}_{i}}{m_i}
\end{aligned}
\end{equation}

That gives us the missing $\Delta \mathbf{p}_{total}$

\begin{equation}
\begin{aligned}
\label{none}
\Delta \mathbf{p}_{total,i}=-\sum_k P \mu_{ik} A_{ik} \mathbf{n}_{ik}
\end{aligned}
\end{equation}


This is source-owned, geometric, and GLSL-friendly. It is not automatically globally momentum-conserving unless every source computes its matching response consistently, but it removes the current bad average and gives a concrete area-to-momentum rule.







\subsection{Force}

\begin{equation}
\begin{aligned}
\label{frclaw}
F(\alpha) = Q\alpha^2
\end{aligned}
\end{equation}

Equation \ref{frclaw}, is used to predict the zero-velocity depth: 
\begin{equation}
\begin{aligned}
\label{zerodepth}
\alpha_0 =
\left(\frac{3\mu v_{in}^2}{2Q}\right)^{1/3}
\end{aligned}
\end{equation}


\subsection{Multicontact rule}

The multi-particle rule should is stated as:

The change in velocity is proportional to the total collision overlap area. 
\begin{equation}
\begin{aligned}
\label{mc}
\Delta v_k \propto A_k
\end{aligned}
\end{equation}

The contributing weight of each collision partner is the ratio of target overlap to total overlap.

\begin{equation}
\begin{aligned}
\label{mc}
w_k = \frac{A_k}{\sum_j A_j}
\end{aligned}
\end{equation}

The change in momentum is equal to the total momentum times the weight.

\begin{equation}
\begin{aligned}
\label{mca}
\Delta \mathbf{p}_k = w_k \Delta \mathbf{p}_{total}
\end{aligned}
\end{equation}
